\documentclass[11pt]{article}
\usepackage[a4paper,margin=1in]{geometry}
\usepackage{amsmath,amssymb,mathtools,bm}
\usepackage{enumitem,booktabs,array}
\usepackage{tcolorbox}
\usepackage{hyperref}
\usepackage[protrusion=true,expansion=false]{microtype}
\hypersetup{colorlinks=true,linkcolor=blue,urlcolor=blue,citecolor=blue}
\setlist[itemize]{topsep=2pt,itemsep=2pt}
\setlist[enumerate]{topsep=2pt,itemsep=2pt}
\newtcolorbox{keybox}{colback=blue!3!white,colframe=blue!40!black,boxrule=0.4pt,arc=1mm}
\newtcolorbox{alertbox}{colback=red!3!white,colframe=red!40!black,boxrule=0.4pt,arc=1mm}
\newtcolorbox{answerbox}{colback=green!3!white,colframe=green!40!black,boxrule=0.4pt,arc=1mm}
\newtcolorbox{drillbox}{colback=yellow!5!white,colframe=orange!60!black,boxrule=0.4pt,arc=1mm}
\newcommand{\EV}{\mathbb{E}}
\newcommand{\Var}{\mathrm{Var}}
\newcommand{\Cov}{\mathrm{Cov}}
\newcommand{\blank}[1]{\underline{\hspace{#1}}}

\title{E07-04: Harford, Stanfield \& Zhang (2019, JFE) \\
\large ``Do Insiders Time Management Buyouts and Freezeouts to Buy Undervalued Targets?'' --- Active Recall Workbook}
\author{Empirical Corporate Finance II --- Exam Prep}
\date{\today}
\begin{document}\maketitle

\begin{keybox}
\textbf{Paper in one sentence.} HSZ provide evidence that managers strategically time management buyouts (MBOs) and freezeouts to coincide with private-information-driven troughs in target stock prices, by documenting that MBO/freezeout firms consistently outperform on operating metrics and stock returns post-deal while paying low completion premia --- insiders use information advantages to buy shareholders out cheaply.

\textbf{Placement.} Core evidence on the agency-cost side of going-private transactions; directly relevant to Delaware fiduciary-duty jurisprudence and regulatory treatment of insider-led deals. Pairs with BDS (E07-01) on LBO efficiency.
\end{keybox}

\section*{Round 1: Complete Walk-Through}

\subsection*{1.1 Research question}
MBOs and freezeouts feature insiders (management, controlling shareholders) on both sides of the transaction. Do they time such deals to private-information-based undervaluation, paying shareholders an artificially low premium?

\subsection*{1.2 Setting}
\begin{itemize}
\item MBO: management team + PE backer acquires all public shares.
\item Freezeout: controlling shareholder buys remaining minority.
\item Legal: Delaware ``entire fairness'' standard for insider-led deals.
\end{itemize}

\subsection*{1.3 Data}
\begin{itemize}
\item SDC MBO and freezeout deals 1990--2014, U.S.\ targets.
\item Compustat fundamentals; CRSP returns pre- and post-deal.
\item Comparison arm-length acquisitions with similar target characteristics.
\end{itemize}

\subsection*{1.4 Empirical findings}
\begin{itemize}
\item Deal premia are lower in MBOs/freezeouts than in arms-length acquisitions of similar targets.
\item Pre-deal target stock performance is depressed relative to industry peers.
\item Post-deal (with PE as informed buyer), firms outperform on operating metrics: margins, ROA, growth.
\item IPO-reintroduction or resale transactions later ($\sim$3--5 years) realize high implied returns for insiders.
\item Freezeouts show the most pronounced pattern, consistent with controlling-shareholder informational advantage.
\end{itemize}

\subsection*{1.5 Interpretation}
\begin{itemize}
\item Supports the ``informed-insider timing'' hypothesis.
\item Insiders appear to exploit short-run undervaluation that they are positioned to know about.
\item Complements legal scrutiny under Delaware entire-fairness review.
\end{itemize}

\subsection*{1.6 Policy and governance implications}
\begin{itemize}
\item Independent special committees and majority-of-minority votes should be the norm.
\item Disclosure of insider information advantages should be required.
\item Evidence informs SEC Rule 13e-3 and state-court fairness analyses.
\end{itemize}

\subsection*{1.7 Caveats}
\begin{alertbox}
\begin{itemize}
\item Hard to distinguish timing from genuine operational-improvement opportunity.
\item Selection of targets into MBO/freezeout is endogenous.
\item Post-deal outcomes may reflect PE value-add rather than pre-deal undervaluation.
\end{itemize}
\end{alertbox}

\section*{Round 2: Level 1 Fill-in}
\begin{enumerate}
\item MBO = \blank{2cm} + PE buys public shares; freezeout = controlling \blank{2cm}.
\item Premium pattern: \blank{1cm} in MBO/freezeout vs.\ arm-length.
\item Pre-deal returns: \blank{2cm}.
\item Post-deal: \blank{2cm} operating performance.
\item Legal standard: Delaware ``\blank{2cm} fairness.''
\end{enumerate}

\section*{Round 3: Level 2 Prompts}
\begin{enumerate}
\item Distinguish MBOs from freezeouts.
\item Describe the empirical pattern (premia, pre-deal returns, post-deal operations).
\item Explain the informed-insider hypothesis.
\item Relate HSZ to Delaware fiduciary doctrine.
\item Evaluate endogeneity concerns.
\end{enumerate}

\section*{Round 4: Minimal Scaffolding}
From a blank page: (i) setting, (ii) data, (iii) findings, (iv) theory, (v) policy.

\section*{Round 5: Blank Page}
Essay: insider timing in going-private transactions.

\vspace{6pt}
\begin{drillbox}
\textbf{Speed Drill.} (1) MBO vs.\ freezeout. (2) Premium. (3) Pre-deal returns. (4) Post-deal performance. (5) Legal standard.
\end{drillbox}

\section*{Quick Reference \& Answer Key}
\begin{answerbox}
\begin{itemize}
\item \textbf{Deals.} MBO \& freezeout by insiders.
\item \textbf{Premia.} Lower than arms-length.
\item \textbf{Post.} Targets outperform later.
\item \textbf{Lesson.} Informed-insider timing.
\end{itemize}
\end{answerbox}

\begin{keybox}
\textbf{30-second exam pitch.} Harford, Stanfield \& Zhang (2019) study MBOs (management + PE) and freezeouts (controlling shareholders buying remaining minority) to test whether insiders time such deals to private-information-based undervaluation. Using SDC deals 1990--2014 and comparing to arm's-length acquisitions with similar targets, they find deal premia are significantly lower in MBO/freezeout transactions, pre-deal target stock performance is depressed relative to peers, and post-deal firms outperform on margins and growth, with particularly strong patterns in freezeouts. The results support an informed-insider-timing hypothesis: insiders exploit short-run undervaluation and PE expertise to buy shareholders out cheaply. The paper is central to Delaware ``entire fairness'' jurisprudence and to SEC Rule 13e-3, arguing for independent special committees, majority-of-minority votes, and stricter disclosure to protect outside shareholders from insider-led going-private transactions.
\end{keybox}

\end{document}
